% !TEX program = pdflatex
% Rémi Zallot -- Curriculum Vitae
%
% Compile with:
%   pdflatex cv.tex && biber cv && pdflatex cv.tex && pdflatex cv.tex
%
% Requires only a standard TeX Live / MacTeX install (tgheros, biblatex,
% biber, fontawesome5 are all bundled) -- no OS-level font installation
% needed, since tgheros loads TeX Gyre Heros through LaTeX's own font
% system rather than searching installed system fonts.

\documentclass[10pt,a4paper]{article}

\usepackage[utf8]{inputenc}
\usepackage[T1]{fontenc}
\usepackage{tgheros}
\renewcommand{\familydefault}{\sfdefault}

\usepackage[margin=2cm]{geometry}
\usepackage{xcolor}
\definecolor{accent}{HTML}{0F766E}
\definecolor{darktext}{HTML}{1F2937}
\definecolor{graytext}{HTML}{6B7280}
\color{darktext}

\usepackage{titlesec}
\usepackage{enumitem}
\usepackage{fontawesome5}
\usepackage{ragged2e}
\usepackage{microtype}

\usepackage[
  colorlinks=true,
  linkcolor=accent,
  urlcolor=accent,
  citecolor=accent,
  pdftitle={Rémi Zallot - Curriculum Vitae},
  pdfauthor={Rémi Zallot},
  pdfsubject={Curriculum Vitae},
  pdfkeywords={biochemistry, bioinformatics, genomics, enzyme function, curriculum vitae},
  pdfusetitle
]{hyperref}

\usepackage[backend=biber,style=numeric-comp,sorting=ydnt,defernumbers=true,maxnames=11,minnames=10]{biblatex}
\addbibresource{publications.bib}

\titleformat{\section}
  {\Large\bfseries\color{accent}}{}{0em}{}
  [{\color{accent}\titlerule[0.8pt]}]
\titlespacing{\section}{0pt}{1.3em}{0.7em}

\titleformat{\subsection}
  {\bfseries\color{darktext}}{}{0em}{}
\titlespacing{\subsection}{0pt}{0.6em}{0.2em}

\setlist[itemize]{leftmargin=1.1em, itemsep=0.15em, topsep=0.15em, parsep=0pt}

\newcommand{\cvdate}[1]{\textcolor{graytext}{\small #1}}
\newcommand{\cventry}[4]{%
  \noindent\textbf{#1}\hfill\cvdate{#2}\\
  \textit{#3}\\
  #4\par\vspace{0.4em}
}

\pagestyle{empty}

\begin{document}

% ---------------------------------------------------------------- Header
\begin{center}
  {\Huge\bfseries Dr Rémi Zallot}\\[0.35em]
  {\large Senior Lecturer (Associate Professor) in Biochemistry}\\[0.15em]
  {\large Department of Life Sciences, Manchester Metropolitan University, UK}\\[0.6em]
  \small
  \faEnvelope\ \href{mailto:remi.zallot@mmu.ac.uk}{remi.zallot@mmu.ac.uk} \quad\textbullet\quad
  \faMapMarker*\ Manchester, United Kingdom \quad\textbullet\quad
  \faGlobe\ \href{https://remizallot.github.io}{remizallot.github.io}\\[0.25em]
  ORCID: \href{https://orcid.org/0000-0002-7317-1578}{0000-0002-7317-1578} \quad\textbullet\quad
  \faGithubSquare\ \href{https://github.com/remizallot}{remizallot} \quad\textbullet\quad
  \faGraduationCap\ \href{https://scholar.google.com/citations?user=sMUlExUAAAAJ}{Google Scholar}
\end{center}

% ---------------------------------------------------------------- Academic positions
\section*{Academic Positions}

\cventry{Senior Lecturer (Associate Professor) in Biochemistry}{Aug 2024 -- present}
  {Department of Life Sciences, Manchester Metropolitan University, UK}{}

\cventry{Lecturer (Assistant Professor) in Biochemistry}{Jun 2023 -- Aug 2024}
  {Department of Life Sciences, Manchester Metropolitan University, UK}{}

\cventry{Research Associate}{-- 2023}
  {Manchester Institute of Biotechnology (MIB), The University of Manchester, UK}{}

\cventry{Marie Skłodowska-Curie Individual Fellow}{Oct 2019 -- Mar 2022}
  {Institute of Life Science, Swansea University Medical School, UK}
  {``deCrYPtion'' project (ID 839116), decrypting \textit{Mycobacterium} cytochrome P450 physiological functions using comparative genomics. Supervisor: Prof. Steven L. Kelly.}

\cventry{Enzyme Function Initiative team}{2016 -- 2019}
  {Carl R. Woese Institute for Genomic Biology, University of Illinois at Urbana-Champaign, USA}
  {Enhanced the EFI genomic enzymology web tools; characterised genes from the human gut microbiome; trained biochemists in comparative genomics approaches.}

\cventry{Postdoctoral Researcher}{-- 2016}
  {Microbiology and Cell Science Department, University of Florida, USA}
  {Characterised genes involved in B vitamin metabolism and queuosine tRNA modification in plants and microbes.}

% ---------------------------------------------------------------- Education
\section*{Education}

\cventry{PhD, Biochemistry / Plant Biology}{2011}
  {Laboratoire de Biogenèse Membranaire}
  {Thesis: \textit{Identification and characterization of a lipase expressed during Arabidopsis thaliana reserves hydrolysis}.}

\cventry{Master's, Plant Biology and Biotechnology}{}
  {Université de Bordeaux, France}{}

\cventry{Licence, Biochemistry}{}
  {Université de Bordeaux, France}{}

% ---------------------------------------------------------------- Research interests
\section*{Research Interests}

My aim is to advance our fundamental understanding of microbes to enhance human health. I focus on characterising and understanding the functions of genes from microbes, particularly pathogens, to leverage this knowledge for developing new drugs and biotechnological applications. My approach is multidisciplinary, combining bioinformatics --- including comparative genomics and genomic enzymology --- with experimental techniques from microbiology, genetics, and biochemistry.

% ---------------------------------------------------------------- Funding
\section*{Grants and Funding}

\cventry{The Academy of Medical Sciences -- Springboard Award (SBF009\textbackslash1016)}{Aug 2024 -- Jul 2026}
  {£125,000}
  {STEROSPRING: Characterising an Uncharted Lanosterol Catabolic Pathway in Non-Tuberculous Mycobacterial (NTM) Species.}

\cventry{MMU Research and Knowledge Exchange -- Research Accelerator Grant}{Jan -- Dec 2023}
  {£7,000}
  {Deciphering an Uncharted Sterol Catabolic Pathway and Its Relevance in Non-Tuberculous Mycobacterial Infection.}

\cventry{Manchester-Melbourne-Toronto (MMT) Research Funds}{Apr 2023}
  {£6,960 -- Grant P122678}
  {Missing Links in Global Sulphur Metabolism. Instructor for Genomic Enzymology and Enzyme Function Initiative training, Melbourne, Australia.}

\cventry{Marie Skłodowska-Curie Individual Fellowship}{Oct 2019 -- Mar 2022}
  {€212,933 -- European Commission, Horizon 2020, Grant 839116}{}

\cventry{Higher Education Investment and Recovery Fund}{Aug 2021 -- Feb 2022}
  {£33,954 -- Higher Education Funding Council for Wales (HEFCW)}{}

\cventry{Inception Programme (Investissement d'Avenir) -- Travel Grant}{Jun 2019}
  {€3,000 -- Grant ANR-16-CONV-0005}
  {Workshop Instructor, ``From Genes to Function: A Survival Guide to Web-Based Tools for Microbiologists,'' Institut Pasteur, Paris.}

\cventry{Sêr Cymru II COFUND Fellowship}{Jul 2019}
  {£173,000 -- Grant SU190}
  {Offer declined due to concurrent MSCA Individual Fellowship.}

\cventry{Wales Marie Skłodowska-Curie Individual Fellowship Summer School}{May 2018}
  {£2,000 travel grant}{}

\cventry{Allocation de recherche (PhD Research Allowance Award)}{Sep 2008 -- Nov 2011}
  {€90,000}
  {French Ministry of Higher Education and Research (Ministère de l'Enseignement Supérieur et de la Recherche).}

% ---------------------------------------------------------------- Publications
\section*{Publications}
\nocite{*}
\printbibliography[heading=none]

% ---------------------------------------------------------------- Teaching
\section*{Teaching}

\cventry{Module Leader, Biomedical Cell Biology (BCB)}{2023 -- present}
  {Manchester Metropolitan University, Department of Life Sciences}
  {Final-year optional module for students on the accredited Biomedical Science programme.}

\cventry{Medical Microbiology}{2023 -- present}
  {Manchester Metropolitan University, Department of Life Sciences}
  {Teaching contributor; supervises final-year student projects.}

% ---------------------------------------------------------------- Talks
\section*{Talks and Presentations}

\begin{itemize}[leftmargin=1.1em]
\item \textbf{Summer 2026 meeting of the Acid Fast Club} --- Conference attended. Liverpool School of Tropical Medicine, Liverpool, UK. July 2026.
\item \textbf{Department Summer research symposium Seminar} --- Oral presentation. Manchester Metropolitan University, Department of Life Sciences, Manchester, UK. April 2026.
\item \textbf{Invited Seminar} --- Invited talk. Institute of Medical Research, Northwestern Polytechnical University, Xi'an, China. April 2026.
\item \textbf{Invited Seminar} --- Invited talk. Shanghai Jiao Tong University Medical School, Shanghai, China. April 2026.
\item \textbf{Tuberculosis: Understanding the Disease Across Scales} --- Poster. Keystone Symposium, Cape Town, South Africa. March 2026.
\item \textbf{Winter Acid Fast Club meeting} --- Conference attended. Acid Fast Club, London, UK. January 2026.
\item \textbf{Leicester 2025 summer session of the Acid Fast Club meeting} --- Oral presentation. Acid Fast Club, Leicester, UK. July 2025.
\item \textbf{Invited talk} --- National University of Singapore, Singapore. April 22--26, 2024.
\item \textbf{Poster presentation} --- Gordon Research Conference on Enzymes, Coenzymes and Metabolic Pathways, Waterville Valley, NH, USA. July 16--21, 2023.
\item \textbf{Invited talk and EFI workshop} --- National University of Singapore, Singapore. November 5--10, 2023.
\item \textbf{Invited talk} --- Chinese Academy of Sciences, Institute of Microbiology, Beijing, China. December 21, 2023.
\item \textbf{Invited talk} --- Northwest University, Xi'an, China. December 24, 2023.
\item \textbf{Understanding antimicrobial resistance requires characterising unknown genes} --- Invited talk. Manchester Institute of Biotechnology Fellowship days, The University of Manchester, Manchester, UK. January 26, 2022.
\item \textbf{Comparative genomics and genomic enzymology approaches for functional assignment} --- Poster (selected for a short oral presentation). Manchester Institute of Biotechnology Annual Science Day, Manchester, UK. 2022.
\item \textbf{Similarity Networks and Genome Neighborhood Networks for function assignment} --- Invited talk. School of Chemistry, Cardiff University, Cardiff, UK. April 3, 2021.
\item \textbf{Genomic enzymology web tools for functional assignment} --- Invited talk. Institut Pasteur, Centre of Bioinformatics, Biostatistics and Integrative Biology, Paris, France. June 20, 2019.
\item \textbf{QueI and its associated activator QueJ are non-orthologous replacement of the nitrile reductase QueF in the queuosine biosynthesis pathway} --- Poster. Gordon Research Seminar and Conference, Enzymes, Coenzymes and Metabolic Pathways, Waterville Valley, NH, USA. 2018.
\item \textbf{QueI and its associated activator QueJ are non-orthologous replacement of the nitrile reductase QueF in the queuosine biosynthesis pathway} --- Poster. Frontiers in Metallobiochemistry Symposium and Bioinorganic Workshop, State College, PA, USA. 2018.
\item \textbf{Comparative genomics for the identification of unknown genes, enzymes and transporters involved in Queuosine metabolism} --- Invited talk. University of Illinois at Urbana-Champaign, USA. August 5, 2016.
\item \textbf{Production of molecules of interest from plants in Yeast, can comparative genomics help?} --- Invited talk. Université François Rabelais, Tours, France. May 6, 2016.
\item \textbf{Is Streptomyces secondary metabolism regulation involving lipid rafts?} --- Invited talk. Université Paris 11, Orsay, France. May 17, 2016.
\item \textbf{Identification of missing enzymes and transporters involved in the synthesis and salvage of Queuosine by comparative genomics} --- Poster. EMBO Conference, Ribosome Structure and Function, Strasbourg, France. 2016.
\item \textbf{Identification of missing enzymes and transporters involved in the synthesis and salvage of Queuosine by comparative genomics} --- Poster. DOE JGI User meeting, Genomics of Energy and Environment Meeting, Walnut Creek, CA, USA. 2016.
\item \textbf{Identification of genes involved in Queuosine Salvage} --- Poster. Gordon Research Conference, RNA editing, Lucca (Barga), Italy. 2015.
\item \textbf{Salvage of the thiamin pyrimidine moiety by plant TenA proteins lacking an active site cysteine} --- Invited talk. Plant Biology conference, Portland, OR, USA. July 12--16, 2014.
\item \textbf{Identification of maize and Arabidopsis genes for the salvage of thiamin breakdown products} --- Poster. Gordon Research Conference, Plant Metabolic Engineering, Waterville Valley, NH, USA. 2013.
\item \textbf{Identification of a potential queuine salvage gene in eukaryotes} --- Poster. XXIV tRNA Conference, Olmué, Chile. 2012.
\item \textbf{Identification of a triacylglycerol-lipase from Brassica napus by functional proteomics} --- Poster (selected for a short oral presentation). 5th European Symposium on Plant Lipids, Gdansk, Poland. 2011.
\end{itemize}

% ---------------------------------------------------------------- Service
\section*{Professional Service and Memberships}

\begin{itemize}
\item Member, Acid Fast Club (since July 2026)
\item Expert Evaluator, European Commission Horizon Europe programme
\item Member, UKRI Talent Peer Review College (PRC)
\item Bacterial lead, Genetically Modified Organisms Health and Safety Committee, Manchester Metropolitan University
\item Carbon Literacy Champion, Department of Life Sciences, Manchester Metropolitan University
\item Member, Microbiology Society (since 2019)
\item Member, NTM Network UK, Basic Science \& Microbiology interest group
\item Fellow of the Higher Education Academy (FHEA)
\end{itemize}

\end{document}
